\documentclass[../main.tex]{subfiles}
\usepackage[utf8]{inputenc}
\usepackage[T1]{fontenc}
\usepackage[indonesian]{babel}

\begin{document}

\chapter*{Pemetaan Silabus terhadap Bab Buku}
\addcontentsline{toc}{chapter}{Pemetaan Silabus terhadap Bab Buku}

Buku ini disusun untuk mendukung rencana pembelajaran semester (RPS) selama 16 pertemuan. Berikut adalah pemetaan antara minggu pertemuan di silabus dengan bab yang tersedia di dalam buku ini:

\begin{longtable}{|c|l|l|}
\hline
\textbf{Minggu} & \textbf{Topik Silabus} & \textbf{Bab dalam Buku} \\
\hline
1 & Pengantar \& Urgensi Kriptografi & Bab 1: Pengantar Kriptografi \\
\hline
2 & Layanan Keamanan \& Serangan & Bab 2: Layanan Keamanan Informasi \\
\hline
3 & Matematika Dasar Kriptografi & Bab 3: Landasan Matematika Dasar \\
\hline
4 & Cipher Klasik I & Bab 4: Kriptografi Klasik (Bagian 1) \\
\hline
5 & Cipher Klasik II & Bab 4: Kriptografi Klasik (Bagian 2) \\
\hline
6 & Stream Cipher & Bab 5: Kriptografi Modern Simetris: Stream Cipher \\
\hline
7 & Block Cipher \& DES & Bab 6: Kriptografi Modern Simetris: Block Cipher \\
\hline
8 & \textbf{UTS} & --- \\
\hline
9 & AES \& Mode Operasi & Bab 7: Analisis Algoritma Block Cipher: DES \& AES \\
\hline
10 & Kunci Publik \& RSA & Bab 8: Analisis Algoritma Block Cipher: AES \\
\hline
11 & DH, ElGamal, ECC & Bab 9: Block Cipher Lainnya \\
\hline
12 & Hash \& MAC & Bab 10: Kriptografi Kunci-Publik: RSA \\
\hline
13 & Tanda Tangan Digital \& PKI & Bab 11: Pertukaran Kunci: Diffie-Hellman \\
\hline
14 & SSL/TLS \& Studi Kasus & Bab 12: Kriptografi Kunci-Publik: ElGamal \\
\hline
15 & Review Materi & Bab 13: Kriptografi Kunci-Publik: Knapsack \\
\hline
16 & \textbf{UAS} & Bab 14: Kriptografi Kurva Eliptik (ECC) \\
\hline
\end{longtable}
\textit{Catatan: Penomoran bab telah disesuaikan untuk mengakomodasi kedalaman materi yang lebih detail dibandingkan pembagian minggu sederhana.}

\end{document}
